\documentclass[11pt]{article}
\usepackage[top=1in, bottom=1in, left=1in, right=1in]{geometry}

\usepackage{amsmath}
\usepackage{amssymb}

\usepackage{graphicx}

\title{Howard Math 273, HW\# 3, \\ {\normalsize Fall 2026; Instructor: Sam Hopkins; Due: Friday, October 30}}
\date{}

\begin{document}

\maketitle

\thispagestyle{empty}

\vspace{-1.8cm}

\begin{enumerate}

\item Recall the Stirling numbers of the $2$nd kind are $S(n,k) := \#$ set partitions of $\{1,2,\ldots,n\}$ into $k$ (non-empty) blocks. We saw $F_k(x) := \sum_{n\geq 0}S(n,k)x^n$ satisfies $F_k(x) = \frac{x^k}{(1-x)(1-2x)\cdots(1-kx)}$ for any $k \geq 1$. Find the partial fraction decomposition of $F_k(x)$, i.e., find the numbers $a_j \in \mathbb{Q}$ for which $F_k(x) = a_0 + \frac{a_1}{(1-x)} + \frac{a_2}{(1-2x)} + \cdots + \frac{a_k}{(1-kx)}$. Conclude that $S(n,k) = \sum_{j=0}^{k} a_j \cdot j^n$. \\
{\bf Hint}: Clear denominators, and then plug in $x=\frac{1}{1}, \frac{1}{2}, \frac{1}{3}, \cdots, \frac{1}{k}$ and finally $x=0$. \\
{\bf Bonus (don't have to do)}: Give another proof of the formula $S(n,k) = \sum_{j=0}^{k} a_j \cdot j^n$ using the Principle of Inclusion-Exclusion, by counting ways to ``place balls into boxes.''

\item Recall that a \emph{rooted, labeled tree on $n$ vertices} is a tree (i.e., connected, acyclic graph) with vertex set $[n] := \{1,\ldots,n\}$, where one vertex is designated the root. A \emph{leaf} of a rooted tree is a vertex $v$ for which either: $v$ is not the root and has degree one; or $v$ is the root and has degree zero. (In the usual genealogical picture, a leaf is a vertex with no children.) \\
Let $t_{n,m}$ denote the number of rooted, labeled trees on $n$ vertices with exactly $m$ leaves.
\begin{enumerate}
\item Form the generating function $T(x,y) := \sum_{n \geq 1} \sum_{m=1}^{n} t_{n,m}\, y^m \, \frac{x^n}{n!}$. Explain why we have that $T(x,y) = x(e^{T(x,y)} - 1 + y)$.
\item Use the functional equation from (a) together with Lagrange inversion to show that $t_{n,m} = \frac{n!}{m!} S(n-1,n-m)$, where $S(n,k)$ denotes the Stirling numbers of the 2nd kind (as in problem 1). You may use the fact, which we showed in class, that the Stirling numbers of the 2nd kind have exponential generating function $\sum_{n\geq 0} S(n,k) \frac{x^n}{n!} = \frac{(e^x-1)^k}{k!}$.
\end{enumerate}
{\bf Difficult bonus (don't have to do)}: Give a \emph{bijective} proof that $t_{n,m} = \frac{n!}{m!} S(n-1,n-m)$.

\item \emph{(Stanley, EC1, \#2.25(a))} Let $f_i(m,n)$ be the number of $m\times n$ matrices of $0$'s and $1$'s, with a total of $i$ $1$'s, and with at least one $1$ in each row and column. Use the Principle of Inclusion-Exclusion to show that
\[ \sum_{i \geq 0} f_i(m,n) t^i = \sum_{k=0}^{n}(-1)^k\binom{n}{k}( (1+t)^{n-k}-1)^m.\]

\item \emph{(Stanley, EC1, \#2.25(b))} With $f_i(m,n)$ as in the previous problem, show that
\[ \sum_{m,n \geq 0} \left(\sum_{i \geq 0} f_i(m,n) t^i \right) \frac{x^m}{m!} \frac{y^n}{n!} = e^{-x-y} \cdot \sum_{m,n \geq 0}  (1+t)^{mn} \frac{x^m}{m!} \frac{y^n}{n!} .\]
{\bf Hint}: You can start with the formula from the previous problem, and then do some algebraic manipulations. Alternatively, you can use the theory of exponential generating functions.

\renewcommand{\binom}{\genfrac{(}{)}{0pt}{}}
\DeclareRobustCommand{\qbinom}{\genfrac{\lbrack}{\rbrack}{0pt}{}}
\end{enumerate}


\end{document}
